\documentclass[11pt,a4paper]{article}
\usepackage[a4paper,margin=14mm,top=12mm,bottom=14mm]{geometry}
\usepackage{fontspec}
\setmainfont{Liberation Sans}
\usepackage{polyglossia}
\setdefaultlanguage{french}
\usepackage{amsmath}
\usepackage{xcolor}
\usepackage{graphicx}
\usepackage[most]{tcolorbox}
\usepackage{tabularx,booktabs,array,multirow}
\usepackage{enumitem}
\usepackage{fancyhdr}
\usepackage{lastpage}
\usepackage[hidelinks]{hyperref}
\definecolor{uimm}{HTML}{005677}
\definecolor{cyan}{HTML}{13A2DD}
\definecolor{fond}{HTML}{E6F2FA}
\definecolor{ok}{HTML}{2E7D32}
\definecolor{ko}{HTML}{B3261E}
\definecolor{violet}{HTML}{7B1FA2}
\setlength{\parindent}{0pt}
\setlength{\parskip}{3pt}
\setlist{nosep}
\setlist[itemize]{leftmargin=13pt,label=\textcolor{cyan}{\textbullet}}
\newcommand{\titre}[1]{\vspace{5pt}{\color{uimm}\large\bfseries #1}\par\vspace{-3pt}{\color{cyan}\rule{\linewidth}{1.2pt}}\par\vspace{0pt}}
\tcbset{boite/.style={colback=fond,colframe=uimm,boxrule=0.6pt,arc=2pt,left=6pt,right=6pt,top=4pt,bottom=4pt},
        defi/.style={colback=violet!6,colframe=violet,boxrule=0.8pt,arc=2pt,left=6pt,right=6pt,top=4pt,bottom=4pt,
                     title={\bfseries Défi},coltitle=white,colbacktitle=violet,fonttitle=\small}}
\newcommand{\bandeau}[3]{%
\begin{tcolorbox}[colback=uimm,colframe=uimm,arc=2pt,left=4pt,right=8pt,top=4pt,bottom=4pt]
\begin{minipage}[c]{18mm}\setlength{\fboxsep}{1.2pt}\fcolorbox{white}{white}{\includegraphics[height=16.5mm]{logo_uimm.jpg}}\end{minipage}\hspace{4mm}%
\begin{minipage}[c]{\dimexpr\linewidth-22mm\relax}\color{white}
{\small #1}\\[1pt]{\LARGE\bfseries #2}\\[2pt]{\large #3}\hfill{\small CQPM Technicien qualité}
\end{minipage}
\end{tcolorbox}}
\newcommand{\reponse}[1][1]{\par\vspace{2pt}{\color{gray!60}\foreach \i in {1,...,#1}{\rule{\linewidth}{0.3pt}\par\vspace{9pt}}}}
\usepackage{pgffor}
\pagestyle{fancy}\fancyhf{}
\renewcommand{\headrulewidth}{0pt}
\newcommand{\pied}[1]{\fancyfoot[L]{\footnotesize\color{uimm}#1}\fancyfoot[R]{\footnotesize\color{uimm}Page \thepage{} sur \pageref{LastPage}}}

\usepackage{tikz}
\usepackage{longtable}
\renewcommand{\bandeau}[3]{%
\begin{tcolorbox}[colback=uimm,colframe=uimm,arc=2pt,left=4pt,right=8pt,top=4pt,bottom=4pt]
\begin{minipage}[c]{18mm}\setlength{\fboxsep}{1.2pt}\fcolorbox{white}{white}{\includegraphics[height=16.5mm]{logo_uimm.jpg}}\end{minipage}\hspace{4mm}%
\begin{minipage}[c]{\dimexpr\linewidth-22mm\relax}\color{white}
{\small #1}\\[1pt]{\LARGE\bfseries #2}\\[2pt]{\large #3}
\end{minipage}
\end{tcolorbox}}
% en-tête d'une fiche : code, titre, question, bloc, durée, organisation, niveau
\newcommand{\entete}[7]{%
\begin{tcolorbox}[colback=uimm,colframe=uimm,arc=2pt,left=4pt,right=6pt,top=4pt,bottom=4pt]
\begin{minipage}[c]{18mm}\setlength{\fboxsep}{1.2pt}\fcolorbox{white}{white}{\includegraphics[height=16.5mm]{logo_uimm.jpg}}\end{minipage}\hspace{3mm}%
\begin{minipage}[c]{\dimexpr\linewidth-58mm\relax}\color{white}\raggedright
{\small Base SPC · #4}\\[1pt]{\Large\bfseries #1 · #2}\\[2pt]{\itshape #3}
\end{minipage}\hfill
\begin{minipage}[c]{33mm}\color{white}\small\raggedleft
Durée : \textbf{#5}\\ \textbf{#6}\\ Niveau #7 sur 3
\end{minipage}
\end{tcolorbox}}
\newcommand{\tagE}{{\setlength{\fboxsep}{1.5pt}\colorbox{cyan}{\color{white}\scriptsize\bfseries Ellistat}}}
\newcommand{\tagW}{{\setlength{\fboxsep}{1.5pt}\colorbox{uimm}{\color{white}\scriptsize\bfseries Page web}}}
\tcbset{pred/.style={colback=orange!6,colframe=orange!80!black,boxrule=0.6pt,arc=2pt,left=6pt,right=6pt,top=3pt,bottom=1pt,
                    title={\bfseries Je prédis, avant de commencer},coltitle=white,colbacktitle=orange!80!black,fonttitle=\small},
        retenir/.style={colback=fond,colframe=uimm,boxrule=0.6pt,arc=2pt,left=6pt,right=6pt,top=3pt,bottom=3pt,
                    title={\bfseries À retenir},coltitle=white,colbacktitle=uimm,fonttitle=\small},
        plus/.style={colback=violet!6,colframe=violet,boxrule=0.6pt,arc=2pt,left=6pt,right=6pt,top=3pt,bottom=3pt,
                    title={\bfseries Pour aller plus loin},coltitle=white,colbacktitle=violet,fonttitle=\small},
        act/.style={colback=uimm!10,colframe=uimm,boxrule=0.5pt,arc=2pt,left=5pt,right=5pt,top=2pt,bottom=2pt,before skip=7pt,after skip=3pt},
        verif/.style={colback=orange!6,colframe=orange!80!black,boxrule=0.6pt,arc=2pt,left=6pt,right=6pt,top=3pt,bottom=3pt,
                    title={\bfseries À vérifier sur poste avant la première séance},coltitle=white,colbacktitle=orange!80!black,fonttitle=\small}}
\newcommand{\fichier}[1]{\texttt{\small #1}}
\newcommand{\menu}[1]{\textsf{\textbf{#1}}}
\newcommand{\champ}[1]{\rule{#1}{0.3pt}}

\usepackage{amssymb}
\newcommand{\cb}{$\square$\,}
\newcommand{\etape}[2]{\par\vspace{9pt}{\setlength{\fboxsep}{3pt}\colorbox{uimm}{\color{white}\bfseries\large\ #1\ }}\hspace{4pt}{\color{uimm}\large\bfseries #2}\par\vspace{2pt}}
\newcommand{\motnouveau}[1]{\begin{tcolorbox}[colback=violet!6,colframe=violet,boxrule=0.6pt,arc=2pt,left=6pt,right=6pt,top=3pt,bottom=3pt,title={\bfseries Mot nouveau},coltitle=white,colbacktitle=violet,fonttitle=\small]#1\end{tcolorbox}}
\newcommand{\chiffre}[1]{\begin{tcolorbox}[colback=orange!8,colframe=orange!80!black,boxrule=1pt,arc=3pt,left=8pt,right=8pt,top=5pt,bottom=5pt]#1\end{tcolorbox}}

\pied{Journée du 25/11/2026 · Trame terrain 4 · Avec ou sans cliquet}
\begin{document}
\thispagestyle{fancy}
\bandeau{Journée du 25/11/2026 · TP 4 sur 5 · durée 30 min · en binôme}{Avec ou sans cliquet}{La façon de serrer change-t-elle la mesure ?}
\vspace{-2pt}
\begin{tcolorbox}[boite,left=5pt,right=5pt,top=2pt,bottom=2pt]\small\textbf{Il vous faut.} le même micromètre à cliquet, la même cale de 10 mm, Ellistat. Même opérateur pour les deux séries.\end{tcolorbox}

\etape{1}{Je prédis}
Si je serre à la main, sans le cliquet, mes valeurs seront :\quad \cb plus grandes \quad \cb plus petites \quad \cb plus dispersées \quad \cb pareilles
\etape{2}{Je mesure}
\textbf{Série A, avec le cliquet} : 15 mesures, toujours le même nombre de crans. \textbf{Série B, sans le cliquet} : 15 mesures en serrant le tambour à la main, comme d'habitude, sans forcer. Retirer la cale à chaque mesure. Noter les trois décimales.
\par\smallskip\textbf{Série A (avec cliquet)}
\par\vspace{2pt}{\setlength{\tabcolsep}{2pt}\begin{tabular}{|p{18mm}|p{18mm}|p{18mm}|p{18mm}|p{18mm}|p{18mm}|p{18mm}|p{18mm}|}\hline
\rule{0pt}{8.5mm}{\tiny\color{gray}1} & \rule{0pt}{8.5mm}{\tiny\color{gray}2} & \rule{0pt}{8.5mm}{\tiny\color{gray}3} & \rule{0pt}{8.5mm}{\tiny\color{gray}4} & \rule{0pt}{8.5mm}{\tiny\color{gray}5} & \rule{0pt}{8.5mm}{\tiny\color{gray}6} & \rule{0pt}{8.5mm}{\tiny\color{gray}7} & \rule{0pt}{8.5mm}{\tiny\color{gray}8} \\ \hline
\rule{0pt}{8.5mm}{\tiny\color{gray}9} & \rule{0pt}{8.5mm}{\tiny\color{gray}10} & \rule{0pt}{8.5mm}{\tiny\color{gray}11} & \rule{0pt}{8.5mm}{\tiny\color{gray}12} & \rule{0pt}{8.5mm}{\tiny\color{gray}13} & \rule{0pt}{8.5mm}{\tiny\color{gray}14} & \rule{0pt}{8.5mm}{\tiny\color{gray}15} &  \\ \hline
\end{tabular}}\par
\par\smallskip\textbf{Série B (sans cliquet)}
\par\vspace{2pt}{\setlength{\tabcolsep}{2pt}\begin{tabular}{|p{18mm}|p{18mm}|p{18mm}|p{18mm}|p{18mm}|p{18mm}|p{18mm}|p{18mm}|}\hline
\rule{0pt}{8.5mm}{\tiny\color{gray}1} & \rule{0pt}{8.5mm}{\tiny\color{gray}2} & \rule{0pt}{8.5mm}{\tiny\color{gray}3} & \rule{0pt}{8.5mm}{\tiny\color{gray}4} & \rule{0pt}{8.5mm}{\tiny\color{gray}5} & \rule{0pt}{8.5mm}{\tiny\color{gray}6} & \rule{0pt}{8.5mm}{\tiny\color{gray}7} & \rule{0pt}{8.5mm}{\tiny\color{gray}8} \\ \hline
\rule{0pt}{8.5mm}{\tiny\color{gray}9} & \rule{0pt}{8.5mm}{\tiny\color{gray}10} & \rule{0pt}{8.5mm}{\tiny\color{gray}11} & \rule{0pt}{8.5mm}{\tiny\color{gray}12} & \rule{0pt}{8.5mm}{\tiny\color{gray}13} & \rule{0pt}{8.5mm}{\tiny\color{gray}14} & \rule{0pt}{8.5mm}{\tiny\color{gray}15} &  \\ \hline
\end{tabular}}\par
\etape{3}{Je saisis dans Ellistat}
\begin{enumerate}[label=\textbf{\alph*.},leftmargin=16pt,itemsep=1pt]
\item Dans une grille, deux colonnes nommées \textbf{Cliquet} et \textbf{Main} : taper les 15 valeurs de chaque série.
\item \menu{Stats Desc.}, onglet \menu{Détail}. Cliquer sur le nom de la colonne dans la liste « Échantillons » à gauche.
\end{enumerate}
\etape{4}{Je lis un seul chiffre}
\chiffre{Dans le tableau « PARAMÈTRES », ligne \textbf{Ecart type}, pour chaque série :\par\medskip
Avec cliquet : \champ{28mm}\hfill Sans cliquet : \champ{28mm}\hfill\ }
\etape{5}{Je décide}
La série la plus dispersée est\quad \cb avec cliquet \quad \cb sans cliquet.\par\smallskip
Consigne à afficher au poste (une phrase) : \par\vspace{1pt}\reponse[2]
\motnouveau{\textbf{Écart type} : la taille de la dispersion des mesures. Plus il est petit, plus les valeurs sont groupées. Les deux séries sont mesurées par la même personne : seule la façon de serrer change.}

\end{document}
